\section{Normalized length algebras}
\label{sec:lengths}

We establish the length inputs for Theorem~\ref{thm:complex} and
Lemma~\ref{lem:acyclicity}.  The mixed-characteristic input uses the
perfectoid results of Cai--Lee--Ma--Schwede--Tucker.  We give the
characteristic-$p$ argument in detail, including the passage to the
non-Noetherian perfection.

For a Noetherian local ring $D$ of characteristic $p$ and a primary
ideal of definition $H$, write $H^{[q]}$ for the ideal generated by the
$q$th powers of elements of $H$, where $q=p^n$. Its Hilbert--Kunz
multiplicity is
\[
 e_{\mathrm{HK}}(H,D)
   =\lim_{q=p^n\longrightarrow\infty}
        \frac{\length_D(D/H^{[q]})}{q^{\dim D}}.
\]
The existence of this limit is Monsky's theorem~\cite{Monsky}; see
also \cite[Section~2.2]{Ma}.

\begin{theorem}\label{thm:length-inputs}
Let $(D,\mathfrak d,k)$ be a complete Noetherian local domain of dimension
$h>0$, with perfect residue field $k$ of characteristic $p>0$.
There exist a $D$-algebra $C$ and a function $\lambda$, with values in
$[0,\infty]$, on $C$-modules annihilated by some power of $\mathfrak d$,
with the following properties.
\begin{enumerate}
\item The function is additive on short exact sequences and vanishes on
the zero module.  Direct sums, sums, and directed colimits of modules of
length zero have length zero whenever a fixed power of $\mathfrak d$
annihilates all modules under consideration.
\item For every system of parameters $z_1,\ldots,z_h$ of $D$ and every
integer $a\geq 1$,
\[
 \lambda\bigl(C/(z_1^a,\ldots,z_h^a)C\bigr)
   =a^h e((z_1,\ldots,z_h),D),
 \qquad
 \lambda H_i(z_1^a,\ldots,z_h^a;C)=0\quad(i>0).
\]
In particular, the number
$\mu=\lambda(C/(z_1,\ldots,z_h)C)$ is positive and finite.
\item If $H$ is $\mathfrak d$-primary, then
\[
 \begin{cases}
 \lambda(C/HC)=e_{HK}(H,D),&\operatorname{char}D=p,\\
 \lambda(C/HC)\leq e(H,D),&\operatorname{char}D=0.
 \end{cases}
\]
\end{enumerate}
The same zero-length stability applies to modules killed by a fixed
power of an ideal extended from any $\mathfrak d$-primary ideal of $D$.
\end{theorem}

For $I=(z_1,\ldots,z_h)C$, assertions (1) and (2) supply the
stability, parameter-length, and Koszul hypotheses of
Theorem~\ref{thm:complex}, with $\mu=e((z),D)$.
Assertion (3) then connects the resulting zeroth-homology bound to
ordinary multiplicities: it identifies the quotient length with
Hilbert--Kunz multiplicity in characteristic $p$, and bounds it by
Hilbert--Samuel multiplicity in mixed characteristic.

\subsection{The normalized length formalism}

We use Faltings' normalized length~\cite[Appendix~2]{Faltings},
developed further by Shimomoto~\cite[Section~2]{Shimomoto}, in the
formulation of \cite[Section~2.2]{CLMST}. Let $(A,\mathfrak
m_A,k)$ be one of the complete regular local rings
\[
 A=k[[t_1,\ldots,t_h]]
 \quad\hbox{or}\quad
 A=W(k)[[t_2,\ldots,t_h]],
\]
and in the second case put $t_1=p$.  Choose compatible $p$-power roots
and set
\[
 A_n=A[t_1^{1/p^n},\ldots,t_h^{1/p^n}],
 \qquad \Lambda=\varinjlim_n A_n.
\]
Thus $\Lambda=A_{\mathrm{perf}}$ in characteristic $p$.  In mixed
characteristic $\Lambda$ denotes the uncompleted tower; its $p$-adic
completion is denoted by $A_\infty$.

We consider the abelian category of $\Lambda$-modules annihilated by
$\mathfrak m_A^N$ for some integer $N$ depending on the module.  The
annihilation here is uniform over the whole module.  In both constructions
below we choose $A\subseteq D$ with $\mathfrak m_A\subseteq\mathfrak d$
and evaluate lengths by restriction of scalars.  Thus every $C$-module killed
by a power of $\mathfrak d$ belongs to this category.
A finitely presented
module in this category descends to a finite-length $A_n$-module $M_n$,
and its normalized length is
\begin{equation}\label{eq:normalized-finite-stage}
 \lambda_\infty(M_n\otimes_{A_n}\Lambda)
       =p^{-nh}\length_{A_n}(M_n).
\end{equation}
The use of $A_n$-length instead of $A$-length makes no difference,
because the two rings have the same residue field.  For a finitely
generated module the length is the infimum of the lengths of finitely
presented modules surjecting onto it.  For an arbitrary module it is
then defined by
\begin{equation}\label{eq:normalized-supremum}
 \lambda_\infty(M)=
 \sup\{\lambda_\infty(N):N\subseteq M
                      \text{ is finitely generated over }\Lambda\}.
\end{equation}
These definitions are compatible and give an additive function on
short exact sequences in this category
\cite[Definition~2.2.1 and Propositions~2.2.3, 2.2.6(a)]{CLMST}.
In particular, length is monotone under submodules and quotients.

We record explicitly the consequences for infinite modules that we
will need.

\begin{lemma}\label{lem:normalized-limit}
Suppose that all the modules in the following assertions are
annihilated by a fixed power of $\mathfrak m_A$.
\begin{enumerate}
\item Direct sums, sums, and directed colimits of zero-length modules
have length zero.
\item If $M=\bigcup_n M_n$ is an increasing union of submodules, then
$\lambda_\infty(M)=\sup_n\lambda_\infty(M_n)$.
\item If $M=\varinjlim_n M_n$ and
$\lambda_\infty(M_n)\longrightarrow 0$, then
$\lambda_\infty(M)=0$.  No injectivity assumption on the transition maps
is required in this assertion.
\end{enumerate}
\end{lemma}

\begin{proof}
A finitely generated submodule of a direct sum is contained in a
finite subsum.  If the summands have length zero, additivity gives
length zero for each finite subsum, and
\eqref{eq:normalized-supremum} gives length zero for the direct sum.
Both a sum of submodules and a directed colimit are quotients of
direct sums, proving the first assertion.  The common annihilating
power ensures that those direct sums belong to the length category.

For the second assertion, each finitely generated submodule of $M$ is
contained in some $M_n$.  Apply \eqref{eq:normalized-supremum} and
monotonicity.

For the third assertion, let $N\subseteq M$ be finitely generated over
$\Lambda$.  Representatives of a finite generating set lie in a common
stage $M_{n_0}$.  Consequently $N$ is contained in the image of $M_n$
in $M$ for every $n\geq n_0$.  Monotonicity for submodules and quotients
gives
\[
 0\leq\lambda_\infty(N)\leq\lambda_\infty(M_n)
       \qquad(n\geq n_0).
\]
The right side tends to zero.  Now take the supremum over $N$.
\end{proof}

\subsection{Characteristic \texorpdfstring{$p$}{p}}

Assume first that $D$ has characteristic $p$.  The Cohen--Gabber
normalization theorem~\cite[Theorem~1.1]{KuranoShimomoto} provides
a coefficient field $k\subseteq D$ and
a system of parameters $t_1,\ldots,t_h$ such that
\[
 A=k[[t_1,\ldots,t_h]]\subseteq D
\]
is module-finite and generically separable; this is also the setup of
\cite[Section~3]{CLMST}.  Work inside a fixed algebraic closure of
$\operatorname{Frac}(D)$ and put
\[
 A'=A_{\mathrm{perf}},\qquad C=D_{\mathrm{perf}},\qquad
 K=\operatorname{Frac}(A),\quad L=\operatorname{Frac}(D),\quad
 K'=K_{\mathrm{perf}}=\operatorname{Frac}(A').
\]
We apply $\lambda_\infty$ to $C$-modules by restriction of scalars to
$A'$.  To compute these lengths, we approximate $C$ by algebras
obtained from finite root stages of $D$.

For every $q=p^n$, define
\[
 B_q=A'\otimes_{A^{1/q}}D^{1/q}.
\]
The natural map from this tensor product to $C$ is injective.  Indeed,
$D^{1/q}$ is torsion-free over $A^{1/q}$, and $A'$ is flat over
$A^{1/q}$: the intervening extensions of regular power-series rings
are free.  Thus $B_q$ injects into
\[
 K'\otimes_{K^{1/q}}L^{1/q}.
\]
The latter ring is a field, because the finite separable extension
$L^{1/q}/K^{1/q}$ is linearly disjoint from the purely inseparable
extension $K'/K^{1/q}$.  This identifies $B_q$ with the subring
$A'D^{1/q}$ of $C$.  These subrings increase with $q$ and exhaust $C$:
each contains $D^{1/q}$, and every finite collection of their
coefficients lies in a sufficiently high root extension of $D$.

We give a self-contained form of the finite-stage comparison in
\cite[Section~3 and Theorem~3.0.7]{CLMST}, beginning with a trace-dual
proof of the conductor bound from \cite[Lemma~2.3]{PolstraTucker}.
Although $B_q$ injects into $C$, passing to a primary quotient can
create a kernel.  The conductor bound below controls this kernel by
giving elements that multiply $C$ into $B_q$.

\begin{lemma}\label{lem:perfection-conductor}
There is an element $0\neq g\in A$ such that
\[
 gC\subseteq B_1,
 \qquad g^{1/q}C\subseteq B_q\quad(q=p^n).
\]
\end{lemma}

\begin{proof}
The field $E=LK'$ is a finite separable extension of the perfect field
$K'$, and is therefore perfect.  Since it contains $L$ and is contained
in $L_{\mathrm{perf}}$, it equals $L_{\mathrm{perf}}$.
The ring $A'$ is normal: any integral equation and the finitely many
fractions involved in an element of $\operatorname{Frac}(A')$ descend
to some normal ring $A^{1/q}$.

Choose a $K$-basis $b_1,\ldots,b_r$ of $L$ with $b_j\in D$.
Its trace-dual basis $\beta_1,\ldots,\beta_r$ exists by separability.
Since $(A\setminus\{0\})^{-1}D=L$, there is $0\neq g\in A$ with
$g\beta_j\in D$ for every $j$.  The same two bases are trace-dual
after extending scalars to $K'$.

Every $c\in C$ is integral over $A'$, as is $g\beta_j$.
The trace $\operatorname{Tr}_{E/K'}(g\beta_jc)$ is therefore integral
over $A'$ and lies in $K'$, so normality puts it in $A'$.
The trace-dual expansion now reads
\[
 gc=\sum_{j=1}^r
       \operatorname{Tr}_{E/K'}(g\beta_jc)b_j\in A'D=B_1.
\]
Taking $q$-th roots inside the perfect ring $C$ gives the second
inclusion, because $(A'D)^{1/q}=A'D^{1/q}=B_q$.
\end{proof}

We next calculate quotient lengths. Let $H$ be $\mathfrak d$-primary and put
\[
 M_q=B_q/HB_q,\qquad
 N_q=\operatorname{im}(M_q\longrightarrow C/HC).
\]
Finite-stage normalization gives
\begin{equation}\label{eq:perfection-stage-quotient}
 \lambda_\infty(M_q)=q^{-h}\length_D(D/H^{[q]}).
\end{equation}
To see this directly, $M_q$ is the flat base extension of
$D^{1/q}/HD^{1/q}$ from $A^{1/q}$ to $A'$.  The $q$-th power
isomorphism $D^{1/q}\longrightarrow D$ identifies that quotient with
$D/H^{[q]}$, preserving module length.  Length over $A^{1/q}$ equals
length over $D^{1/q}$ because their residue fields agree.

The kernel of $M_q\longrightarrow N_q$ is annihilated by $g^{1/q}$.
In fact, an element in $B_q\cap HC$ is a finite sum of elements of
$H$ times elements of $C$, and Lemma~\ref{lem:perfection-conductor}
puts its product with $g^{1/q}$ in $HB_q$.
For an endomorphism of a module of finite normalized length,
additivity shows that kernel and cokernel have the same normalized
length.  Apply this to multiplication by $g^{1/q}$ on $M_q$ to obtain
\begin{align}
 \lambda_\infty\ker(M_q\longrightarrow N_q)
 &\leq\lambda_\infty(0:_{M_q}g^{1/q})\notag\\
 &=\lambda_\infty(M_q/g^{1/q}M_q)\notag\\
 &=q^{-h}\length_D D/(H^{[q]},g).
 \label{eq:perfection-kernel-bound}
\end{align}
This last quantity is $O(q^{-1})$.  Indeed, choose
$\mathfrak d^a\subseteq H$, and let $t$ be a number of generators of
$H$.  The pigeonhole containment
\[
 H^{t(q-1)+1}\subseteq H^{[q]}
\]
shows that $H^{[q]}$ contains an ordinary power of $\mathfrak d$ whose
exponent is $O(q)$.  If $g$ is a nonunit, the domain hypothesis gives
$\dim D/(g)\leq h-1$, so Hilbert--Samuel growth on $D/(g)$ gives
$\length_D D/(H^{[q]},g)=O(q^{h-1})$.  If $g$ is a unit, this length
is zero.

Combining \eqref{eq:perfection-stage-quotient} and
\eqref{eq:perfection-kernel-bound}, we obtain
\[
 q^{-h}\length_D D/H^{[q]}-O(q^{-1})
 \leq\lambda_\infty(N_q)
 \leq q^{-h}\length_D D/H^{[q]}.
\]
The submodules $N_q$ increase and exhaust $C/HC$.  They are all killed
by $H$ and hence by one fixed power of $\mathfrak m_A$.  The
Hilbert--Kunz limit and Lemma~\ref{lem:normalized-limit} therefore give
\begin{equation}\label{eq:perfection-hk}
 \lambda_\infty(C/HC)=e_{HK}(H,D).
\end{equation}
The quotient lengths are now identified.  To obtain the remaining
parameter assertions of Theorem~\ref{thm:length-inputs}, we first show
that positive parameter Koszul homology has length zero; the Koszul
Euler characteristic will then identify the parameter quotient length
with Hilbert--Samuel multiplicity.  Fix a system of
parameters $z_1,\ldots,z_h$ of $D$.  Flatness at each stage and the
$q$-th power isomorphism give
\begin{equation}\label{eq:perfection-koszul-stage}
 \lambda_\infty H_i(z_1,\ldots,z_h;B_q)
   =q^{-h}\length_D H_i(z_1^q,\ldots,z_h^q;D).
\end{equation}
These are finite lengths: the finite-stage Koszul homology is a
finitely generated module annihilated by a primary parameter ideal.
The complex $K(z_1,\ldots,z_h;D)$ consists of finite free modules in
degrees $0,\ldots,h$, and its homology has finite length.  The
Frobenius homology estimate of Roberts~\cite{Roberts}, in the form
\cite[Theorem~2.4]{Ma}, thus gives
\[
 \lim_{q=p^n\to\infty}
 q^{-h}\length_D H_i(z_1^q,\ldots,z_h^q;D)=0
 \qquad(i>0).
\]
This theorem requires a complete characteristic-$p$ local ring and a
finite-free complex of length equal to its dimension with finite-length
homology, all of which have just been checked.

Filtered colimits are exact for modules, so
\[
 H_i(z_1,\ldots,z_h;C)
   =\varinjlim_q H_i(z_1,\ldots,z_h;B_q).
\]
Each module here is annihilated by $(z_1,\ldots,z_h)$ and therefore
by a fixed power of $\mathfrak m_A$.  The third assertion of
Lemma~\ref{lem:normalized-limit}, applied to
\eqref{eq:perfection-koszul-stage}, proves
\begin{equation}\label{eq:perfection-koszul-zero}
 \lambda_\infty H_i(z_1,\ldots,z_h;C)=0\qquad(i>0).
\end{equation}
The same argument applies to every fixed parameter list
$z_1^a,\ldots,z_h^a$.

Finally, the Koszul Euler characteristic formula for parameter
ideals~\cite[Tag~0AZZ]{Stacks} gives
\[
 \sum_{i=0}^h(-1)^i\length_D
             H_i(z_1^q,\ldots,z_h^q;D)
   =e((z_1^q,\ldots,z_h^q),D)
   =q^h e((z_1,\ldots,z_h),D).
\]
Divide by $q^h$ and use the positive-homology limits.  The degree-zero
term tends to the Hilbert--Kunz multiplicity of the parameter ideal.
Consequently
\[
 e_{HK}((z_1,\ldots,z_h),D)=e((z_1,\ldots,z_h),D).
\]
Together with \eqref{eq:perfection-hk} and the parameter multiplicity
power rule, it gives the quotient formulas in
Theorem~\ref{thm:length-inputs} in characteristic $p$.

\paragraph{Finite-length modules.}
For the Dutta-multiplicity consequence in
Section~\ref{sec:charp}, we also need the same comparison for a
finite-length module, since zeroth homology need not be cyclic.
For a $D$-module $M$, write
\[
 \Phi_D^n(M)=D\otimes_{D,\varphi_D^n}M
\]
for scalar extension along the $n$th Frobenius endomorphism
$\varphi_D^n:D\to D$.

\begin{lemma}\label{lem:perfection-finite-length}
In the characteristic-$p$ setup above, let $M$ be a finite-length
$D$-module.  If the limit on the right exists, then
\[
 \lambda_\infty(C\otimes_D M)
   =\lim_{q=p^n\longrightarrow\infty}
           q^{-h}\length_D\bigl(\Phi_D^n(M)\bigr).
\]
\end{lemma}

\begin{proof}
The assertion is immediate when $M=0$.  Otherwise put
$J=\Ann_D M$, a $\mathfrak d$-primary ideal, and let $b$ be a number
of generators of $M$.  Define
\[
 M_q=B_q\otimes_D M,\qquad
 N_q=\im(M_q\longrightarrow C\otimes_D M).
\]
The module $M_q$ is obtained by flat scalar extension from
$A^{1/q}$ to $A'$ of
$D^{1/q}\otimes_D M$.  Under the $q$th power isomorphism
$D^{1/q}\cong D$, the latter module corresponds to $\Phi_D^n(M)$.
It is a finite-length module over the Noetherian ring $A^{1/q}$,
so finite-stage normalization gives
\begin{equation}\label{eq:perfection-module-stage}
 \lambda_\infty(M_q)=q^{-h}\length_D\bigl(\Phi_D^n(M)\bigr).
\end{equation}
Here the lengths over $A^{1/q}$ and $D^{1/q}$ agree because their
residue fields agree.

By Lemma~\ref{lem:perfection-conductor}, multiplication by $g^{1/q}$
defines a $D$-linear map $C\to B_q$ whose composite
$B_q\hookrightarrow C\to B_q$ is multiplication by $g^{1/q}$.
After tensoring with $M$, this shows that
$K_q=\ker(M_q\to N_q)$ is killed by $g^{1/q}$.  The finite normalized
length in \eqref{eq:perfection-module-stage} permits the same kernel--cokernel
calculation as in \eqref{eq:perfection-kernel-bound}.  Since $M$ is a
quotient of $(D/J)^b$, it gives
\begin{align*}
 \lambda_\infty(K_q)
 &\leq \lambda_\infty(M_q/g^{1/q}M_q)\\
 &\leq b\,\lambda_\infty\bigl(B_q/(J,g^{1/q})B_q\bigr)\\
 &=bq^{-h}\length_D D/(J^{[q]},g)=O(q^{-1}).
\end{align*}
The last estimate is the Hilbert--Samuel growth argument from
\eqref{eq:perfection-kernel-bound}, with $J$ in place of $H$:
$J^{[q]}$ contains a power of
$\mathfrak d$ with exponent $O(q)$, and $\dim D/(g)\leq h-1$ if
$g$ is a nonunit.  If $g$ is a unit, the error is zero.

The submodules $N_q$ increase to $C\otimes_D M$ and are all killed
by $J$, hence by one fixed power of $\mathfrak m_A$.  Additivity,
\eqref{eq:perfection-module-stage}, and
Lemma~\ref{lem:normalized-limit} now give the assertion by taking
the limit.
\end{proof}

\subsection{Mixed characteristic}

Assume now that $D$ has characteristic zero.  Complete $p$ to a system
of parameters $p,t_2,\ldots,t_h$ of $D$.  Cohen structure gives a
module-finite inclusion, with the same residue field,
\[
 A=W(k)[[t_2,\ldots,t_h]]\subseteq D.
\]
It is generically separable because the fraction fields have
characteristic zero.  Let $A_\infty$ be the $p$-adic completion of the
tower above, and set
\[
 C=D^{A_\infty}_{\mathrm{perfd}}
   :=(D\otimes_A A_\infty)_{\mathrm{perfd}}.
\]
The construction uses the perfectoidization of Bhatt--Scholze
\cite[Section~8 and Theorem~10.11]{BhattScholze}; $C$ is the
perfectoid algebra in \cite[Notation~4.0.1]{CLMST}.
Use its normalized length
$\lambda_\infty$, or equivalently restrict scalars to the uncompleted
tower $\Lambda$ when evaluating bounded
$\mathfrak m_A$-power-torsion modules.  Indeed, $\Lambda$ is
$p$-torsion-free, so $\Lambda/p^N\cong A_\infty/p^N$ for every $N$.
A module annihilated by $\mathfrak m_A^N$ is annihilated by $p^N$;
its submodules and finite generating sets are therefore the same
over the two length bases.

For clarity, the exact perfectoid consequences used here are as
follows.  Choose $0\neq g\in A$ such that $A[1/g]\to D[1/g]$ is
finite etale, and write $(g)_{\mathrm{perfd}}$ for the almost ideal in
\cite[Notation~4.0.1]{CLMST}.  Corollary~4.0.4 of that paper says that
the positive Koszul homology on $C$ of every system of parameters of
$A_n\otimes_A D$ is annihilated by $(g)_{\mathrm{perfd}}$.
Taking $n=0$ covers every system of parameters of $D$, not just the
chosen parameters of $A$.
Proposition~4.0.10 identifies the normalized lengths of
$(g)_{\mathrm{perfd}}$-almost-isomorphic
$\mathfrak m_A$-power-torsion modules.  In particular the Koszul
homology just described has length zero.

Definition~4.0.8 and Corollary~4.0.18 of \cite{CLMST} give, for every
$\mathfrak d$-primary ideal $H$,
\[
 \lambda_\infty(C/HC)=e_{\mathrm{perfd}}(H,D)\leq e(H,D),
\]
with equality when $H$ is a parameter ideal.  The equality for
arbitrary parameter ideals is also the $n=0$ case of
Proposition~4.0.14 there.  Applying it to
$(z_1^a,\ldots,z_h^a)$ and using the ordinary parameter multiplicity
power rule yields
\[
 \lambda_\infty\bigl(C/(z_1^a,\ldots,z_h^a)C\bigr)
       =a^h e((z_1,\ldots,z_h),D).
\]
All these statements apply to the same algebra $C$ and the same
normalized length, while $z$ ranges over the systems of parameters
of $D$.

We finish by checking the category and stability assertions in both
characteristics.  A module killed by a power of $\mathfrak d$ is killed
by a power of $\mathfrak m_A$, so the length above is defined on it and
is additive there.  If $H$ is $\mathfrak d$-primary and $H^N M=0$,
choose $r$ with $\mathfrak d^r\subseteq H$; then
$\mathfrak m_A^{rN}M=0$.  Thus a common power of $H$ supplies the common
annihilating power required by Lemma~\ref{lem:normalized-limit}.
The value $e((z),D)$ is positive and finite.  These observations prove
all the remaining assertions of Theorem~\ref{thm:length-inputs}.
